%% FT_QuestionPool.tex
%% ─────────────────────────────────────────────────────────────────────────────
%% CONTROL QUESTIONS — Question Pool
%% Foreign Trade and International Organisations, Jan Suchanek, Hochschule Bremen, Summer Semester 2026
%%
%% Covers: Ch. 1 Introduction · Ch. 2 The World Economy and National Accounts
%%         Ch. 3 Freight and Transport · Ch. 4 Terms of Trade and the Effects of Trade
%%         Ch. 5 Trade Theory · Ch. 6 Presentations: Countries and Markets
%%         Ch. 7 The State, Taxes and Tariffs · Ch. 8 Subsidies in an Open Economy
%%         Ch. 9 Multinational Enterprises · Ch. 10 International Law and Governance
%%         Ch. 11 Monetary Theory · Ch. 12 Macroeconomics in the Open Economy
%%
%% Each chapter provides: 3 Multiple Choice · 1 True/False · 1 Open Question
%% ─────────────────────────────────────────────────────────────────────────────
\documentclass[paper=a5]{scrartcl}

\usepackage{xcolor}
\usepackage{hsbcommon}
\usepackage[hidelinks]{hyperref}
\usepackage[utf8]{inputenc}
\usepackage{setspace}
\usepackage{pgfplots}
\pgfplotsset{compat=1.18}

%% ── Load the question-bank package FIRST ────────────────────────────────────
\usepackage{hsbquestions}

%% ── TOGGLE — uncomment EXACTLY ONE line ─────────────────────────────────────
\showanswerstrue               % ANSWER KEY — reveals answers + solutions
% \showanswersfalse            % STUDENT VERSION — hides answers (default)

\usepackage[headsepline]{scrlayer-scrpage}
\automark{section}
\pagestyle{scrheadings}
\ihead{\leftmark}
\chead{}
\ohead{Control Questions}

\setlength{\parskip}{1ex}

\title{Foreign Trade and International Organisations}
\subtitle{Control Questions — Full Question Pool}
\author{Sonnet 4.6 / Claude AI using the Script of Jan \noun{Suchanek}}

\begin{document}
\setcounter{page}{247}
\onehalfspacing

\maketitle

\clearpage



%% ═════════════════════════════════════════════════════════════════════════════
\section{Chapter 1 — The World Economy and National Accounts}
%% ═════════════════════════════════════════════════════════════════════════════

\begin{question}[type=mc]{The World as a Closed Economy}
The world economy is described as the only ``closed economy''. What does
this imply?
\begin{choices}
  \choice        Global trade is regulated by a single world government.
  \correctchoice No external trading partner exists, so every import is
                 simultaneously someone else's export, and the global current
                 account is always zero.
  \choice        Global GDP cannot grow because there is no new demand.
  \choice        All currencies must be pegged to gold to prevent imbalances.
\end{choices}
\begin{solution}
  A closed economy has no rest-of-world sector. For the planet as a whole,
  every buyer has a seller and every import is someone else's export, so the
  \emph{global} current account must net to zero. This does not mean that
  individual countries are in balance. The world economy is also \emph{not}
  a zero-sum game: global production and wealth can grow over time through
  innovation and specialisation.
\end{solution}
\end{question}

\begin{question}[type=mc]{Orders of Magnitude}
Why does the script stress the importance of memorising a handful of
``standard candle'' economic figures (such as your home country's GDP)?
\begin{choices}
  \choice        To pass the written exam without additional study.
  \correctchoice To have reference points that give large economic numbers
                 meaningful perspective, making it easier to judge whether a
                 quoted figure is plausible.
  \choice        Because the WTO requires all traders to know national GDPs.
  \choice        To be able to convert between currencies without a calculator.
\end{choices}
\begin{solution}
  When dealing with large numbers --- the cost of an oil tanker, the daily
  grain trade volume, the GDP of a commodity-dependent country --- it helps to
  anchor them against known figures. If you know Germany's GDP is roughly
  \$4 trillion, you can quickly judge whether a stated figure for a smaller
  economy is plausible. The script suggests treating these reference figures
  as mental ``standard candles''.
\end{solution}
\end{question}

\begin{question}[type=mc]{Fundamental Balance of the World Economy}
Which of the following is a direct consequence of the world economy being
a closed system?
\begin{choices}
  \choice        Rich countries always grow faster than poor ones.
  \choice        Trade is only beneficial when globally balanced.
  \choice        Currency reserves are always constant worldwide.
  \correctchoice Every unit of global exports equals a unit of global imports;
                 no country can run a surplus without another running a deficit.
\end{choices}
\begin{solution}
  Because the world has no external counterpart, every export transaction is
  simultaneously an import transaction for the receiving country. In aggregate,
  global exports equal global imports identically. An individual country can
  accumulate a surplus, but this is exactly mirrored by deficits elsewhere.
  This accounting identity holds regardless of exchange rate regimes or trade
  policies.
\end{solution}
\end{question}

\begin{question}[type=tf]{Global Trade is a Zero-Sum Game}
Because every exporter's gain corresponds to an importer's payment, global
trade is a zero-sum game in which total wealth cannot increase.
\truefalse[false]
\begin{solution}
  \textbf{False.} Trade is \emph{not} zero-sum. Voluntary exchange occurs
  because both parties expect to gain --- each values what they receive more
  than what they give up. Additionally, specialisation according to comparative
  advantage raises total output, so the global production frontier expands.
  The world economy grows through trade, innovation, and division of labour,
  not merely through redistribution of a fixed stock of wealth.
\end{solution}
\end{question}

\begin{question}[type=open]{Why Do ``Orders of Magnitude'' Matter in Foreign Trade
  Analysis?}
The script advises students to develop a sense of scale for large economic
figures. Explain why this skill matters when analysing foreign trade data,
and give two concrete examples of how it can expose misleading claims.

\answerlines[6]

\begin{solution}
  Trade data involves very large numbers (billions of USD, millions of tonnes)
  that are easy to misrepresent or misinterpret without reference points.
  A sense of scale allows the analyst to detect implausibly large or small
  claims and to contextualise statistics properly.

  \textbf{Example 1:} A news headline claiming that a country ``earned
  \$500 billion from coffee exports'' is immediately suspect if you know that
  the entire global coffee market is worth roughly \$100 billion per year.

  \textbf{Example 2:} A company reporting that it ``controls 5\% of the
  global copper market'' can be evaluated only if you know that annual global
  copper production is around 22 million tonnes --- 5\% would be roughly
  1.1 million tonnes, a significant but plausible quantity for a major miner.

  In both cases, a mental ``standard candle'' --- total market size, a
  country's GDP, a commodity's annual production volume --- allows rapid
  plausibility checking.
\end{solution}
\end{question}


%% ═════════════════════════════════════════════════════════════════════════════
\section{Chapter 2 — Freight and Transport}
%% ═════════════════════════════════════════════════════════════════════════════

\begin{question}[type=mc]{Incoterms}
Incoterms are defined by the International Chamber of Commerce. Their
primary purpose in a trade contract is to specify:
\begin{choices}
  \choice        The currency in which payment must be made.
  \choice        The applicable import tariff rate for the good.
  \correctchoice Which party is responsible for which stage of the delivery
                 chain, including transport, insurance, and the transfer of
                 risk.
  \choice        The quality standards the good must meet at the border.
\end{choices}
\begin{solution}
  Incoterms are standardised three-letter codes (e.g.\ FOB, CIF, EXW, DDP)
  that define the point at which responsibility and risk transfer from seller
  to buyer. They specify who arranges and pays for inland transport, loading,
  ocean freight, insurance, and customs clearance in the importing country.
  Because ownership and risk transfer are among the most legally significant
  moments in any trade contract, Incoterms reduce disputes and transaction
  costs enormously.
\end{solution}
\end{question}

\begin{question}[type=mc]{Contango and Backwardation}
In commodity futures markets, ``contango'' describes a situation where:
\begin{choices}
  \choice        The spot price is higher than the futures price, indicating
                 excess supply.
  \correctchoice The futures price is higher than the spot price, reflecting
                 storage costs, insurance, and the cost of carry.
  \choice        Prices collapse because sellers outnumber buyers.
  \choice        A commodity cannot be traded on an exchange.
\end{choices}
\begin{solution}
  Under normal conditions, \textbf{contango} holds: buying a commodity now
  and storing it costs money (warehouse fees, insurance, financing), so the
  futures price exceeds the spot price. \textbf{Backwardation} is the
  opposite: the spot price exceeds the futures price, which can occur when
  there is an immediate shortage or very high convenience yield. Understanding
  the term structure of futures prices is essential for hedging trade risk.
\end{solution}
\end{question}

\begin{question}[type=mc]{Value Chain: Upstream vs.\ Downstream}
In the commodity value chain, which of the following activities is
\emph{most upstream}?
\begin{choices}
  \choice        Branding and retail distribution to consumers.
  \choice        Assembly of manufactured consumer goods.
  \choice        Refining a raw material into an intermediate product.
  \correctchoice Mining or harvesting the raw material.
\end{choices}
\begin{solution}
  \textbf{Upstream} activities are closest to the resource: mining, drilling,
  harvesting. Moving downstream, raw materials are refined, then used to
  produce intermediate goods, which are assembled into manufactured wares,
  branded, and distributed through wholesale and retail channels.
  Value is added (and costs incurred) at every stage. Most of the final
  consumer price typically reflects downstream activities --- branding,
  marketing, distribution --- not the raw material itself.
\end{solution}
\end{question}

\begin{question}[type=tf]{Physical and Financial Fulfilment Are Independent}
In an international trade contract, the physical delivery of goods and the
financial payment are legally separate events that can occur at different
times and places, and must both be explicitly specified.
\truefalse[true]
\begin{solution}
  \textbf{True.} A trade contract has two sides: \emph{physical fulfilment}
  (the seller transfers the good to the buyer) and \emph{financial fulfilment}
  (the buyer pays the seller). The time, place, and conditions for each must
  be specified in the contract, because they determine when risk and ownership
  transfer, what insurance is needed, and what happens in case of default.
  Incoterms govern the physical side; payment terms (letter of credit, open
  account, etc.) govern the financial side.
\end{solution}
\end{question}

\begin{question}[type=open]{Why Did Technological Progress Historically Drive the
  Expansion of International Trade?}
Using examples from the history of trade outlined in the script, explain how
reductions in transportation costs and financial innovations enabled trade
to expand across longer distances and larger volumes.

\answerlines[6]

\begin{solution}
  \textbf{Transport innovations} directly lowered the cost per tonne-kilometre.
  The discovery of America opened new trade routes; the Suez Canal halved the
  journey from Europe to Asia; the steam boat replaced unpredictable wind
  power; and containerisation in the 20th century standardised loading and
  dramatically cut handling costs. Each innovation made previously unprofitable
  trade routes commercially viable.

  \textbf{Financial innovations} reduced the risk of long-distance trade.
  Paper money and the banking system meant that merchants no longer needed to
  ship precious metals across oceans. Bills of exchange and letters of credit
  enabled payment without physical currency transfer. Insurance products
  covered cargo loss. Futures and forward contracts allowed traders to lock in
  prices and hedge against exchange rate movements during long voyages.

  Together, these innovations made international trade safer, cheaper, and
  faster --- explaining the long-run trend towards greater openness and
  globalisation.
\end{solution}
\end{question}


%% ═════════════════════════════════════════════════════════════════════════════
\section{Chapter 3 — Terms of Trade and the Effects of Trade}
%% ═════════════════════════════════════════════════════════════════════════════

\begin{question}[type=mc]{Winners and Losers from Trade}
When a domestic market opens to free trade and cheaper imports arrive,
which group typically experiences the sharpest welfare loss?
\begin{choices}
  \choice        Consumers of the imported good.
  \correctchoice Domestic producers of the import-competing good.
  \choice        Logistics companies and port operators.
  \choice        Foreign exporters who now supply the market.
\end{choices}
\begin{solution}
  Cheaper imports lower the domestic price, which benefits consumers but
  undercuts domestic producers. Those producers may be forced to reduce
  output, lay off workers, or exit the market entirely. Because they are
  often locally concentrated and well organised, they tend to lobby strongly
  against trade liberalisation --- even though aggregate consumer gains
  typically exceed producer losses. This asymmetry between concentrated
  losers and diffuse winners is a central theme in the political economy
  of trade.
\end{solution}
\end{question}

\begin{question}[type=mc]{Terms of Trade}
The ``terms of trade'' of a country are defined as:
\begin{choices}
  \choice        The tariff rates applied at the border.
  \choice        The Incoterms agreed between buyer and seller.
  \correctchoice The ratio of a country's export price index to its import
                 price index.
  \choice        The legal conditions set by the WTO for bilateral agreements.
\end{choices}
\begin{solution}
  The \textbf{terms of trade} measure the relative price of exports versus
  imports. An improvement (the ratio rises) means each unit of exports buys
  more imports --- a gain in purchasing power. A deterioration hurts countries
  dependent on imports (e.g.\ oil importers after a price spike). Commodity
  exporters are particularly exposed to terms-of-trade shocks, since their
  export revenues depend on world commodity prices they cannot control.
\end{solution}
\end{question}

\begin{question}[type=mc]{Producer Lobbying and Trade Policy}
Why do producers often succeed in obtaining trade protection even when the
aggregate welfare case for free trade is strong?
\begin{choices}
  \choice        International law obliges governments to protect domestic
                 producers.
  \correctchoice Producers are concentrated, well-organised, and face large
                 per-person losses; consumers are diffuse and face small
                 individual losses, so they rarely lobby effectively against
                 protection.
  \choice        Trade theory is silent on the distributional effects of trade.
  \choice        Consumers always benefit from tariffs, so they support
                 protection.
\end{choices}
\begin{solution}
  This is a classic \textbf{political economy} argument. When a factory closes
  in a single town, all affected workers have a strong, visible incentive to
  organise and lobby. The cost of the resulting tariff is spread across
  millions of consumers who each pay a few euros more per year --- too little
  to justify individual political action. The result is systematic
  over-representation of producer interests in trade policy, even when the
  aggregate welfare cost of protection is large.
\end{solution}
\end{question}

\begin{question}[type=tf]{Free Trade Always Benefits Every Citizen}
Free trade is welfare-maximising in the aggregate, which means every citizen
of every trading country is made better off when trade opens.
\truefalse[false]
\begin{solution}
  \textbf{False.} Free trade raises \emph{aggregate} welfare, but the gains
  and losses are distributed unevenly. Consumers and efficient exporters gain;
  producers of import-competing goods lose. In principle, winners could
  compensate losers (Kaldor-Hicks criterion), but this compensation rarely
  occurs automatically. The existence of domestic losers explains why trade
  liberalisation generates political resistance even in countries that benefit
  overall.
\end{solution}
\end{question}

\begin{question}[type=open]{Explain How Trade Creates Both Winners and Losers Within
  a Single Country}
Using the concept of the trade value chain and the supply-demand diagram,
explain why opening a market to free trade benefits some domestic agents
and harms others. What determines who wins and who loses?

\answerlines[7]

\begin{solution}
  When a country opens to trade and the world price of a good is \emph{below}
  the autarky price, imports flow in and the domestic price falls to the world
  level.

  \textbf{Winners:} Domestic consumers gain consumer surplus (the price is
  lower). Traders, logistics companies, and ports benefit from higher trade
  volumes. Domestic producers of goods that are now \emph{exported} (where
  the world price exceeds the autarky price) also gain.

  \textbf{Losers:} Domestic producers of import-competing goods face lower
  prices and may be forced out of the market. Workers in those industries may
  lose jobs. If the affected industry is geographically concentrated (a
  ``company town''), entire communities can be negatively affected.

  \textbf{What determines the split?} The elasticities of supply and demand,
  the share of the good in household budgets, the labour intensity of the
  industry, and the ease of retraining and relocation. The script notes that
  producers are typically better organised and lobby more effectively than
  consumers, which is why protection survives despite aggregate welfare losses.
\end{solution}
\end{question}


%% ═════════════════════════════════════════════════════════════════════════════
\section{Chapter 4 — Trade Theory}
%% ═════════════════════════════════════════════════════════════════════════════

\begin{question}[type=mc]{The Wealth of Nations and Division of Labour}
According to Adam Smith (1776), what is the foundation of a nation's wealth?
\begin{choices}
  \choice        Accumulation of gold and silver reserves.
  \choice        Military strength and control of trade routes.
  \correctchoice The division of labour, which raises productivity through
                 specialisation.
  \choice        High tariff barriers that protect domestic industry.
\end{choices}
\begin{solution}
  Smith's central insight in \emph{The Wealth of Nations} is that \textbf{the
  division of labour} drives productivity. By specialising in narrow tasks,
  workers become more skilled and efficient, and complex goods can be produced
  at lower cost. This principle scales up from individuals to firms to nations:
  countries specialise in what they do relatively best, trade the surplus, and
  import what others produce more cheaply. The result is mutual gain.
\end{solution}
\end{question}

\begin{question}[type=mc]{Comparative Advantage}
Ricardo's principle of comparative advantage states that a country should
export a good when:
\begin{choices}
  \choice        It can produce the good with absolutely fewer inputs than
                 any other country.
  \choice        Its government has declared the good a strategic industry.
  \correctchoice Its opportunity cost of producing that good is lower than
                 its trading partner's, regardless of absolute productivity.
  \choice        World demand for the good is growing faster than world supply.
\end{choices}
\begin{solution}
  \textbf{Comparative advantage} is about \emph{relative} (opportunity) costs,
  not absolute productivity. Even if Country A is more productive than
  Country B in every sector, trade is still mutually beneficial: each country
  specialises where its opportunity cost is lowest. Country A exports the
  good in which its productivity advantage is largest; Country B exports the
  good in which its disadvantage is smallest. Both end up consuming more than
  in autarky.
\end{solution}
\end{question}

\begin{question}[type=mc]{Reasons for Trade}
Which of the following is \emph{not} listed in the script as a reason
why countries trade?
\begin{choices}
  \choice        To obtain goods that require a different climate or resource
                 base.
  \choice        To dispose of surplus production.
  \correctchoice To comply with mandatory WTO membership obligations.
  \choice        To maximise profits through arbitrage and resale.
\end{choices}
\begin{solution}
  The script identifies genuine economic motivations: obtaining goods you
  cannot produce (different climate/resources/know-how), exporting surpluses,
  exploiting comparative advantage, and profit-seeking through arbitrage.
  WTO membership imposes disciplines on trade policy (tariff bindings, MFN),
  but it does not \emph{compel} countries to trade; it constrains how they
  may restrict trade. Trade itself is driven by economic incentives, not
  legal mandates.
\end{solution}
\end{question}

\begin{question}[type=tf]{Dumping Is Always Illegal}
Selling goods on foreign markets at a price below the domestic price or
below the cost of production --- a practice known as dumping --- is
automatically illegal under WTO rules.
\truefalse[false]
\begin{solution}
  \textbf{False.} Dumping is not automatically illegal under WTO rules.
  The WTO Anti-Dumping Agreement permits importing countries to impose
  anti-dumping duties, but only after a formal investigation confirms
  dumping has occurred, domestic industry has been harmed, and the
  two are causally linked. The process is lengthy and contested. Many
  dumping cases are driven by political pressure from domestic industries
  rather than clear economic harm.
\end{solution}
\end{question}

\begin{question}[type=open]{Why Are the Decisions of What to Produce and What to
  Trade ``Intrinsically Linked''?}
The script states that ``we do not just trade things that we happen to have,
but we have things so that we can trade.'' Explain this idea and its
implications for how we understand the gains from trade.

\answerlines[6]

\begin{solution}
  In a simple barter or endowment model, trade is purely about exchanging
  surpluses. But in a modern economy, production decisions are made
  \emph{in anticipation of} trade. A coffee farmer in Ethiopia does not
  grow coffee for personal consumption; they grow it to sell on world
  markets and purchase other goods with the proceeds. A German car
  manufacturer designs and produces vehicles specifically for export.

  This means the gains from trade are not limited to exchanging pre-existing
  surpluses --- trade actively \emph{reshapes} what countries produce. Countries
  develop industries in which they have comparative advantages because world
  markets reward specialisation. Without trade, those industries would not
  exist at their current scale.

  The implication is that the gains from trade are larger and more dynamic
  than a static endowment model suggests: trade encourages learning-by-doing,
  scale economies, and investment in productive capacity aligned with
  comparative advantage.
\end{solution}
\end{question}



%% ═════════════════════════════════════════════════════════════════════════════
\section{Chapter 5 — The State, Taxes and Tariffs}
%% ═════════════════════════════════════════════════════════════════════════════

\begin{question}[type=mc]{Government Expenditure Ratio}
The government expenditure ratio (state quota) is calculated as:
\begin{choices}
  \choice        Tax revenue divided by total government debt.
  \correctchoice Total government spending divided by GDP.
  \choice        The share of the workforce employed by the public sector.
  \choice        Defence spending as a percentage of the national budget.
\end{choices}
\begin{solution}
  The \textbf{state quota} $SQ = G / Y$ measures how large the public
  sector is relative to the economy. It varies widely: France at around 58\%,
  Germany at 49\%, the USA at 36\%, Saudi Arabia at 28\%, Pakistan at 20\%.
  High-income countries tend to have larger state quotas, reflecting more
  extensive welfare states. War efforts and crises (as Ukraine illustrates)
  sharply increase the ratio.
\end{solution}
\end{question}

\begin{question}[type=mc]{Role of the Central Bank}
Which of the following best describes the primary functions of a central bank?
\begin{choices}
  \choice        To compete with commercial banks for retail deposits.
  \correctchoice To maintain price stability, regulate the banking system,
                 manage the money supply via interest rates and open market
                 operations, and act as lender of last resort.
  \choice        To negotiate international trade agreements on behalf of the
                 government.
  \choice        To collect income taxes and distribute social transfers.
\end{choices}
\begin{solution}
  The central bank is a unique public institution that holds the privilege of
  issuing currency. Its main tools include setting the \textbf{interest rate},
  conducting \textbf{open market operations} (buying/selling government bonds),
  setting \textbf{minimum reserve requirements}, and regulating commercial
  banks. Most modern central banks are formally \emph{independent} of
  government to prevent short-term political pressure from compromising
  price stability.
\end{solution}
\end{question}

\begin{question}[type=mc]{Externalities and the Case for Taxation}
A factory produces steel but also emits pollutants that impose health costs
on nearby residents. In the absence of regulation, the factory will:
\begin{choices}
  \choice        Voluntarily reduce pollution to maximise social welfare.
  \correctchoice Over-produce, because it does not bear the external costs
                 of pollution, so private marginal cost is below social
                 marginal cost.
  \choice        Under-produce, because it pays compensation to residents.
  \choice        Produce at the socially optimal level, because markets are
                 efficient.
\end{choices}
\begin{solution}
  \textbf{Externalities} are costs (or benefits) imposed on third parties
  that are not reflected in market prices. When pollution is a negative
  externality, the firm's private marginal cost lies \emph{below} the social
  marginal cost. The firm therefore produces too much relative to the social
  optimum. A \textbf{Pigouvian tax} equal to the marginal external cost can
  internalise the externality and restore the efficient output level.
\end{solution}
\end{question}

\begin{question}[type=tf]{The Statutory and Economic Incidence of a Tariff Always
  Coincide}
The party legally required to pay a tariff --- typically the importer ---
always bears the full economic burden of that tariff.
\truefalse[false]
\begin{solution}
  \textbf{False.} The \emph{statutory} incidence (who writes the cheque)
  differs from the \emph{economic} incidence (who ultimately bears the cost).
  An importer pays the tariff to customs but may pass some or all of it
  forward to consumers through higher prices, or absorb part of it through
  lower margins depending on the elasticities of supply and demand. In a
  small country, a tariff raises the domestic price and the burden falls on
  domestic consumers; foreign exporters are unaffected.
\end{solution}
\end{question}

\begin{question}[type=open]{What Are the Main Types of Government Policy, and How
  Do They Interact in an Open Economy?}
The script distinguishes between fiscal, regulatory, and monetary policy.
Explain each briefly and discuss how they interact when a country is exposed
to international trade.

\answerlines[7]

\begin{solution}
  \textbf{Fiscal policy} involves government spending and borrowing (and
  implicitly, taxation). Expansionary fiscal policy raises aggregate demand;
  in an open economy, some of that demand leaks abroad as imports, reducing
  the domestic multiplier effect.

  \textbf{Regulatory policy} covers rules, standards, tariffs, quotas, and
  obligations imposed on private actors. Trade-related regulations (tariffs,
  non-tariff barriers, product standards) directly affect import and export
  volumes.

  \textbf{Monetary policy} is managed by the (ideally independent) central
  bank via interest rates and open market operations. In an open economy,
  interest rate changes affect capital flows and the exchange rate: a rate
  rise attracts capital inflows, appreciating the currency, which makes
  exports more expensive and imports cheaper.

  \textbf{Interaction:} A fiscal stimulus can be partially offset if the
  central bank raises rates to contain inflation, appreciating the currency
  and reducing net exports. Tariffs that raise net exports have a
  macroeconomic effect similar to fiscal stimulus (rightward shift of IS).
  The three policy types must therefore be considered jointly in an open
  economy context.
\end{solution}
\end{question}


%% ═════════════════════════════════════════════════════════════════════════════
\section{Chapter 6 — Subsidies in an Open Economy}
%% ═════════════════════════════════════════════════════════════════════════════

\begin{question}[type=mc]{Production Subsidy vs.\ Export Subsidy}
A \emph{production subsidy} differs from an \emph{export subsidy} in that:
\begin{choices}
  \choice        A production subsidy is funded by tariff revenue; an export
                 subsidy is not.
  \choice        Only an export subsidy is subject to WTO disciplines.
  \correctchoice A production subsidy lowers costs for both domestic and
                 export sales; an export subsidy specifically targets goods
                 sold abroad and may lower their price below the domestic
                 market price.
  \choice        There is no meaningful economic difference between the two.
\end{choices}
\begin{solution}
  A \textbf{production subsidy} reduces a producer's costs across the board:
  output for both home and foreign markets increases. A \textbf{export subsidy}
  targets only exported goods, driving a wedge between the domestic and
  foreign price. This is particularly distorting because it can displace
  efficient foreign producers and may constitute dumping. The WTO's Agreement
  on Subsidies and Countervailing Measures (SCM) disciplines both, but export
  subsidies face stricter rules.
\end{solution}
\end{question}

\begin{question}[type=mc]{Welfare Effect of a Production Subsidy in an Open Economy}
In a small open economy, a production subsidy to a domestic industry
(at the world price):
\begin{choices}
  \choice        Raises the domestic price above the world price.
  \choice        Leaves consumer welfare unchanged and reduces producer
                 welfare.
  \correctchoice Increases domestic production and reduces imports, while
                 leaving the domestic consumer price unchanged (equal to the
                 world price).
  \choice        Creates a deadweight gain equal to the value of the subsidy.
\end{choices}
\begin{solution}
  Because the small country is a price-taker, the domestic consumer price
  remains at the world price. The subsidy lowers the effective cost to
  domestic producers, so they expand output. Since the market still clears
  at the world price, imports simply fall by the amount of the additional
  domestic production. Consumers are unaffected; producers gain; but the
  government (taxpayers) bears the subsidy cost. The net welfare effect is
  negative: the subsidy cost exceeds the producer surplus gain (deadweight
  loss triangle).
\end{solution}
\end{question}

\begin{question}[type=mc]{Dumping and Export Subsidies}
An export subsidy that causes the export price to fall below the cost of
production may constitute dumping. Why is this considered harmful by
importing countries?
\begin{choices}
  \choice        It raises prices for consumers in the importing country.
  \correctchoice Artificially low prices can drive efficient domestic producers
                 out of business; once competitors are eliminated, the
                 subsidising country may raise prices to monopoly level.
  \choice        It automatically violates GATT Article I (MFN principle).
  \choice        It reduces the subsidising country's GDP.
\end{choices}
\begin{solution}
  The harm from dumping is dynamic: in the short run, consumers in the
  importing country benefit from low prices. But if domestic producers exit
  the market and cannot easily re-enter, the exporter acquires market power
  and may subsequently raise prices. This is the \emph{predatory pricing}
  logic applied to trade. Importing countries may respond with
  \textbf{anti-dumping duties} under WTO rules after a formal injury
  investigation.
\end{solution}
\end{question}

\begin{question}[type=tf]{Subsidies Always Increase Net Social Welfare}
A subsidy transfers resources from taxpayers to a domestic industry, but
because it increases output and reduces imports, it always raises total
net social welfare.
\truefalse[false]
\begin{solution}
  \textbf{False.} A production subsidy creates a \emph{deadweight loss}: the
  subsidy cost paid by taxpayers exceeds the producer surplus gain. The
  additional domestic production is supplied at a cost higher than the world
  price (otherwise domestic producers would have supplied it without the
  subsidy), meaning resources are used less efficiently than if the same
  goods were simply imported. Total welfare in the subsidising country falls.
  Subsidies may be justified on non-economic grounds (strategic industries,
  food security, environmental goals), but the economic efficiency cost is real.
\end{solution}
\end{question}

\begin{question}[type=open]{How Do Export Subsidies Distort International Trade,
  and How Does the WTO Address Them?}
Explain the mechanism by which an export subsidy distorts trade, who gains
and who loses, and what instruments the international trading system provides
to counter the distortion.

\answerlines[7]

\begin{solution}
  An \textbf{export subsidy} artificially lowers the price at which a
  country's producers can sell abroad. This expands exports, displaces
  competing producers in importing countries (who face unfair competition),
  and potentially depresses world prices for the subsidised commodity.

  \textbf{Who gains:} Domestic producers in the subsidising country and
  consumers in importing countries (lower prices). \textbf{Who loses:}
  Taxpayers in the subsidising country (who fund the subsidy); domestic
  producers in importing countries; efficient third-country exporters who
  lose market share.

  \textbf{WTO instruments:}
  \begin{itemize}
    \item The \emph{SCM Agreement} prohibits certain export subsidies (``red
      light'') and allows countervailing duties against injurious subsidies.
    \item The importing country may initiate a \emph{countervailing duty}
      investigation and impose offsetting tariffs if injury is proven.
    \item Disputes can be brought to the \emph{WTO Dispute Settlement Body}.
  \end{itemize}
  Agricultural export subsidies were a major issue in WTO negotiations and
  were finally eliminated (for most products) under the 2015 Nairobi
  Decision.
\end{solution}
\end{question}


%% ═════════════════════════════════════════════════════════════════════════════
\section{Chapter 7 — Multinational Enterprises}
%% ═════════════════════════════════════════════════════════════════════════════

\begin{question}[type=mc]{Herfindahl-Hirschman Index}
The Herfindahl-Hirschman Index (HHI) is calculated as $H = \sum_i a_i^2$,
where $a_i$ is each firm's market share. What does the HHI measure?
\begin{choices}
  \choice        A company's export intensity relative to domestic sales.
  \correctchoice Market concentration: the higher the HHI, the more the
                 market is dominated by a few large firms.
  \choice        The degree of vertical integration in a supply chain.
  \choice        A country's balance of trade in a specific commodity.
\end{choices}
\begin{solution}
  The \textbf{HHI} ranges from near zero (perfect competition with many
  small firms) to 1 (or 10{,}000 if expressed in percentage points), which
  represents pure monopoly. Competition authorities (the EU Commission,
  the US Department of Justice) use it to assess whether mergers create
  problematic market concentration. A post-merger HHI above 2{,}500 (in
  percentage-point form) typically triggers scrutiny.
\end{solution}
\end{question}

\begin{question}[type=mc]{Transfer Pricing}
Transfer pricing is relevant to multinational enterprises because:
\begin{choices}
  \choice        It determines the official exchange rate used for
                 intra-group payments.
  \correctchoice Prices set for transactions between subsidiaries of the
                 same corporate group can be manipulated to shift profits to
                 lower-tax jurisdictions, reducing the group's overall tax
                 bill.
  \choice        It sets the tariff applied when goods move between
                 subsidiaries in different customs unions.
  \choice        It governs the cost of shipping goods between plants.
\end{choices}
\begin{solution}
  \textbf{Transfer pricing} refers to the prices charged on transactions
  within a multinational group (e.g.\ a subsidiary in a high-tax country
  selling an intermediate product to a subsidiary in a low-tax country at
  below-market price). By understating profits in high-tax countries and
  overstating them in low-tax havens, MNEs reduce their global tax burden.
  The OECD's \emph{Base Erosion and Profit Shifting} (BEPS) project and
  country-by-country reporting requirements are designed to counter this.
\end{solution}
\end{question}

\begin{question}[type=mc]{Market Forms and Competition}
Which market structure is characterised by many sellers offering
differentiated products, so that each firm has some price-setting power
but faces competition from close substitutes?
\begin{choices}
  \choice        Perfect competition.
  \choice        Oligopoly.
  \correctchoice Monopolistic competition.
  \choice        Monopsony.
\end{choices}
\begin{solution}
  \textbf{Monopolistic competition} (Chamberlin) combines features of both
  monopoly and competition: many firms, free entry and exit, but each firm
  sells a differentiated product and thus faces a downward-sloping demand
  curve. Examples include restaurants, clothing brands, and petrol stations.
  In the long run, free entry erodes economic profit, but each firm still
  produces below the efficient scale (excess capacity theorem).
\end{solution}
\end{question}

\begin{question}[type=tf]{A Company Listed on a Stock Exchange Is Always Publicly
  Owned}
If a company is listed on a stock exchange, it is owned by the government.
\truefalse[false]
\begin{solution}
  \textbf{False.} A \emph{listed} (publicly traded) company has shares that
  can be bought and sold on a stock exchange by private investors, funds,
  and institutions --- it is \emph{publicly traded} but not necessarily
  \emph{publicly owned}. Government (state) ownership means the government
  holds the shares; this can occur in listed companies (partial state
  ownership) but is a separate question from listing status. Many MNEs are
  fully privately owned and listed; others are state-owned enterprises (SOEs)
  that may or may not be listed.
\end{solution}
\end{question}

\begin{question}[type=open]{Why Is It Important to Look Beyond a Company's Own
  Published Information When Analysing an MNE?}
Describe the limitations of using only a company's website and annual report
as analytical sources. What additional information should a thorough analysis
of a multinational enterprise include, and why?

\answerlines[7]

\begin{solution}
  \textbf{Limitations of company-published information:}
  \begin{itemize}
    \item Annual reports are audited but shaped by management; they
      emphasise strengths and legally required disclosures, not strategic
      weaknesses.
    \item Websites are marketing tools aimed at customers, employees, and
      investors, not independent analysis.
    \item Transfer pricing, tax structures, and lobbying activities are
      rarely highlighted.
  \end{itemize}

  \textbf{Additional sources a thorough analysis should include:}
  \begin{itemize}
    \item \emph{Independent financial analysts} (broker reports, rating
      agencies) for profitability, debt levels, and competitive position.
    \item \emph{Industry reports} from bodies such as the IEA, FAO, ICSG,
      or UNCTAD for market context and commodity fundamentals.
    \item \emph{NGO and investigative journalism reports} for environmental,
      labour, and governance controversies.
    \item \emph{Regulatory filings} (SEC 20-F for US-listed foreign firms,
      EU competition decisions) for legally required disclosures.
    \item \emph{Competitors' reports} for an industry-wide perspective.
  \end{itemize}

  Only by triangulating these sources can the analyst form a balanced
  picture of the company's actual performance, risks, and strategic position.
\end{solution}
\end{question}


%% ═════════════════════════════════════════════════════════════════════════════
\section{Chapter 8 — International Law and Governance}
%% ═════════════════════════════════════════════════════════════════════════════

\begin{question}[type=mc]{Why International Law Is Different from National Law}
International law is described as fundamentally different from domestic
national law. What is the key reason?
\begin{choices}
  \choice        International law applies only to corporations, not to states.
  \correctchoice There is no world state and no supranational enforcement
                 authority; all international cooperation is voluntary and
                 relies on consent between sovereign states.
  \choice        International law is written by the UN Security Council and
                 is fully binding on all nations.
  \choice        International law has no written form and relies solely on
                 historical custom.
\end{choices}
\begin{solution}
  The absence of a \textbf{world state} is the defining feature. Domestic law
  can be enforced by courts and police; international law can only be enforced
  through diplomacy, retaliation, trade sanctions, or (in rare cases) UN
  Security Council action. Even the WTO Dispute Settlement Body can only
  authorise retaliation --- it cannot compel compliance. Hugo Grotius
  (1583--1645), the father of international law, recognised this and built his
  framework on the idea of states voluntarily accepting norms of natural law.
\end{solution}
\end{question}

\begin{question}[type=mc]{The Most Favoured Nation Principle}
Article I of GATT requires WTO members to apply the Most Favoured Nation
(MFN) principle. This means:
\begin{choices}
  \choice        The largest trading nation automatically receives the lowest
                 tariff rate.
  \choice        Developing countries are exempt from all import duties.
  \correctchoice A WTO member may not charge any trading partner a tariff
                 higher than the lowest tariff it charges to any other member,
                 with the only permitted exception being membership of a
                 recognised customs union or free trade area.
  \choice        Members must reduce tariffs to zero within ten years of
                 joining.
\end{choices}
\begin{solution}
  \textbf{MFN} is the foundational non-discrimination principle of the
  multilateral trading system: if you give one country a tariff concession,
  you must give it to all WTO members. The only exception is for \emph{free
  trade areas and customs unions} (GATT Article XXIV): members of such
  blocs may give each other preferential treatment without extending it to
  all WTO members. This exception has been the main driver of the explosion
  of Preferential Trade Agreements (PTAs) since the 1990s.
\end{solution}
\end{question}

\begin{question}[type=mc]{Degrees of Economic Integration}
Preferential Trade Agreements vary in depth. Ranking from shallowest to
deepest integration, which order is correct?
\begin{choices}
  \choice        Customs union → free trade agreement → common market →
                 economic union → currency union.
  \choice        Economic union → currency union → common market → customs
                 union → free trade agreement.
  \correctchoice Free trade agreement → customs union → common market →
                 economic union → currency union.
  \choice        Currency union → economic union → common market → free trade
                 agreement → customs union.
\end{choices}
\begin{solution}
  The standard taxonomy of integration (Balassa, 1961) runs: \\
  \textbf{FTA} (no tariffs among members, separate external tariffs) →
  \textbf{Customs Union} (common external tariff) →
  \textbf{Common Market} (free movement of goods, services, capital, labour) →
  \textbf{Economic Union} (common regulatory/policy framework) →
  \textbf{Currency Union} (single currency).
  The EU has progressed to currency union for Eurozone members, while USMCA
  remains a free trade agreement.
\end{solution}
\end{question}

\begin{question}[type=tf]{The WTO Replaced GATT, Which No Longer Exists}
The World Trade Organization replaced the General Agreement on Tariffs and
Trade (GATT) in 1995, and GATT is no longer legally in force.
\truefalse[false]
\begin{solution}
  \textbf{False.} GATT is still legally in force. The WTO incorporated
  GATT 1994 as one of its founding agreements, and the original GATT 1947
  continues to have legal effect through transitional provisions. The
  WTO extended the framework to cover services (GATS) and intellectual
  property (TRIPS), which GATT did not address. The WTO is thus best
  understood as \emph{built on top of} GATT, not as its replacement.
\end{solution}
\end{question}

\begin{question}[type=open]{Why Has the MFN Principle Paradoxically Led to a
  Proliferation of Preferential Trade Agreements?}
Explain the logic of the MFN principle and why --- rather than promoting
uniform global free trade --- it has incentivised the formation of regional
trade blocs.

\answerlines[6]

\begin{solution}
  The \textbf{MFN principle} requires that any tariff concession granted to
  one WTO member be extended to all others. This makes bilateral deals that
  favour only one partner illegal --- unless the two parties form a recognised
  free trade area or customs union under GATT Article XXIV.

  The \textbf{paradox:} Countries wishing to give preferential market access
  to close partners (for geopolitical or economic reasons) cannot do so
  bilaterally under MFN. The only legal route is to form a full FTA or
  customs union. This creates a strong incentive to establish such blocs,
  which proliferated dramatically from the 1990s onwards (EU enlargement,
  NAFTA/USMCA, ASEAN, Mercosur, CETA, JEFTA, etc.).

  The result is a complex ``spaghetti bowl'' of overlapping PTAs with
  different rules, rather than the clean multilateral liberalisation GATT
  intended. The \textbf{ratchet principle} --- tariffs can only fall, never
  rise --- means these concessions lock in, further complicating the
  multilateral system. WTO members have effectively outsourced much of
  their trade liberalisation to the regional level.
\end{solution}
\end{question}


%% ═════════════════════════════════════════════════════════════════════════════
\section{Chapter 9 — Monetary Theory}
%% ═════════════════════════════════════════════════════════════════════════════

\begin{question}[type=mc]{The Three Costs of Money}
The script identifies three distinct costs associated with holding or
lending money internationally. Which set is correct?
\begin{choices}
  \choice        Principal, commission, and exchange fee.
  \correctchoice The interest rate ($i$), inflation ($\pi$), and the
                 risk of currency devaluation ($\Delta\text{Exp}$).
  \choice        Administration costs, custody fees, and transaction charges.
  \choice        Opportunity cost, sunk cost, and marginal cost.
\end{choices}
\begin{solution}
  When money is lent internationally, the lender faces three risks/costs:
  (1)~the \textbf{interest rate} is the explicit price of credit; (2)~
  \textbf{inflation} erodes the real purchasing power of the repayment;
  (3)~\textbf{currency devaluation} reduces the value of a foreign-currency
  repayment when converted back to the home currency. Understanding all three
  is essential for pricing international bonds, setting exchange rate policy,
  and analysing the Balance of Payments.
\end{solution}
\end{question}

\begin{question}[type=mc]{Gresham's Law}
Gresham's Law, famously illustrated by Isaac Newton's failed attempt to
favour silver over gold in England, states that:
\begin{choices}
  \choice        High interest rates always drive foreign capital out of
                 a country.
  \choice        Governments that devalue their currency face inflation.
  \correctchoice ``Bad money drives out good'': when two currencies of
                 different intrinsic value circulate at the same legal
                 exchange rate, the overvalued currency circulates while the
                 undervalued one is hoarded or exported.
  \choice        Paper money is always preferable to commodity money.
\end{choices}
\begin{solution}
  Newton decreed that gold would be exchanged for silver at a ratio of 1:21
  (up from 1:20), believing this would encourage use of silver. Instead,
  gold was now \emph{undervalued} by the new rate: foreigners imported gold
  (cheap by the official rate) and exported silver. This perfectly illustrates
  \textbf{Gresham's Law}: the artificially ``cheap'' money (silver, now
  overvalued officially) circulates; the money whose intrinsic value exceeds
  its official value (gold) disappears from circulation.
\end{solution}
\end{question}

\begin{question}[type=mc]{Structure of the Balance of Payments}
The Balance of Payments is divided into which main accounts?
\begin{choices}
  \choice        Trade Account and Reserve Account only.
  \correctchoice Current Account (goods, services, primary and secondary
                 income) and Financial \& Capital Account (direct investment,
                 portfolio investment, reserve account).
  \choice        Fiscal Account and Monetary Account.
  \choice        Export Account and Import Account.
\end{choices}
\begin{solution}
  The \textbf{Current Account} records flows of goods (trade balance),
  services, primary factor income (wages, investment returns), and
  secondary income (transfers). The \textbf{Financial \& Capital Account}
  records cross-border capital flows: foreign direct investment, portfolio
  investment (bonds, equities), other investment, and changes in reserve
  assets. By definition, the sum of all accounts is zero: a current account
  deficit must be financed by a capital account surplus.
\end{solution}
\end{question}

\begin{question}[type=tf]{Quantitative Easing Creates ``Something from Nothing''}
When a central bank buys government bonds through open market operations
(quantitative easing), it creates new money without any corresponding
liability, permanently increasing the nation's wealth.
\truefalse[false]
\begin{solution}
  \textbf{False.} Quantitative easing (\textbf{QE}) involves the central
  bank purchasing assets (government bonds, mortgage-backed securities) by
  crediting the sellers' reserve accounts. The central bank's balance sheet
  expands symmetrically --- a new asset (the bond) and a new liability
  (the reserve credit). Wealth is not created from nothing; the central bank
  exchanges one asset for another. The economic effect is to lower long-term
  interest rates and increase money supply, stimulating credit and investment.
  Unwinding QE (quantitative tightening) reverses the process.
\end{solution}
\end{question}

\begin{question}[type=open]{Explain the History and Function of the Banking System
  in Facilitating International Trade}
Trace how the development of paper money, the giro system, and the banking
sector reduced the frictions of long-distance trade. What risks did these
innovations introduce, and how did regulation respond?

\answerlines[7]

\begin{solution}
  \textbf{The friction problem:} Early long-distance trade required shipping
  precious metals in the opposite direction of goods --- dangerous, slow, and
  costly. Merchants needed a way to make and receive payment without physically
  moving gold or silver across oceans.

  \textbf{Innovations:}
  \begin{itemize}
    \item The \emph{bill of exchange} (Italian banking, Florence): a written
      promise to pay, allowing merchants to settle accounts by offsetting
      debits and credits at a central exchange point.
    \item The \emph{giro system} (Hamburg 1619, Nuremberg, Venice): banks
      kept accounts and transferred balances without physical currency.
    \item \emph{Paper money and banknotes} (Bank of England 1694, Banque de
      France 1800): convertible claims on the issuing institution, easier
      to transport than coin.
    \item \emph{Stock exchanges and bond markets} (Bruges 1276, Amsterdam
      1602): allowed risks to be pooled and long-term capital to be raised
      for trade ventures.
  \end{itemize}

  \textbf{Risks introduced:} Banks can create ``book money'' beyond their
  reserves, creating systemic risk if confidence collapses (bank runs).

  \textbf{Regulatory response:} Central banks emerged as supervisors and
  lenders of last resort. Minimum reserve requirements and deposit guarantees
  were introduced to prevent self-fulfilling bank runs.
\end{solution}
\end{question}


%% ═════════════════════════════════════════════════════════════════════════════
\section{Chapter 10 — Macroeconomics in the Open Economy}
%% ═════════════════════════════════════════════════════════════════════════════

\begin{question}[type=mc]{The IS Curve}
In the IS-LM model, the IS curve represents:
\begin{choices}
  \choice        All combinations of interest rate and output where the money
                 market is in equilibrium.
  \correctchoice All combinations of interest rate and output where the real
                 goods market is in equilibrium, i.e.\ where investment equals
                 savings ($I = S$).
  \choice        The relationship between inflation and unemployment (the
                 Phillips curve).
  \choice        The effect of government spending on the exchange rate.
\end{choices}
\begin{solution}
  The \textbf{IS curve} is derived from the equilibrium condition in the
  goods market: investment ($I$) equals savings ($S$). Higher interest rates
  make borrowing expensive, depressing investment and hence output --- so the
  IS curve slopes downward. Fiscal policy (government spending $G$, net
  exports $NX$) shifts the IS curve: a rise in $G$ or in $NX$ (e.g.\ from
  a tariff) shifts IS to the right, raising both output and the interest rate.
\end{solution}
\end{question}

\begin{question}[type=mc]{Crowding Out}
In the IS-LM model, when the government increases spending and shifts the IS
curve to the right, the resulting higher interest rate partially offsets the
stimulus. This phenomenon is called:
\begin{choices}
  \choice        Quantitative easing.
  \correctchoice Crowding out.
  \choice        The Laffer effect.
  \correctchoice Crowding out.
\end{choices}
\begin{solution}
  \textbf{Crowding out} occurs because the fiscal stimulus raises output and
  income, increasing the demand for money (transactions motive). This pushes
  up the interest rate, which reduces private investment. The net increase in
  output is therefore \emph{less than} the initial government spending
  increase. In an open economy, additional crowding out can occur if higher
  interest rates attract capital inflows and appreciate the currency, reducing
  net exports.
\end{solution}
\end{question}

\begin{question}[type=mc]{Effect of Expansionary Monetary Policy in the IS-LM
  Model}
A central bank expands the money supply, shifting the LM curve to the right.
The new IS-LM equilibrium shows:
\begin{choices}
  \choice        Higher interest rates and higher output.
  \choice        Lower interest rates and lower output.
  \correctchoice Lower interest rates and higher output.
  \choice        Higher interest rates and lower output.
\end{choices}
\begin{solution}
  A rightward shift of the \textbf{LM curve} (more money available at each
  income level) lowers the interest rate for a given level of output. Lower
  rates stimulate investment and consumption, raising output. The new
  equilibrium is at a lower interest rate and higher output. However, the
  stimulus is also only partial: higher output raises money demand, partially
  offsetting the initial fall in the interest rate. This is why monetary and
  fiscal policy are sometimes used together.
\end{solution}
\end{question}

\begin{question}[type=tf]{A Tariff in the IS-LM Framework Has No Macroeconomic
  Effect}
Because tariffs are a microeconomic trade policy tool, they have no effect
on the macroeconomic IS or LM curves.
\truefalse[false]
\begin{solution}
  \textbf{False.} A tariff reduces imports, raising net exports
  $NX = EX - IM$. An increase in $NX$ is an increase in aggregate demand,
  which shifts the \textbf{IS curve to the right} --- exactly the same
  macroeconomic effect as an increase in government spending. The script
  explicitly states this. Tariffs therefore have macroeconomic consequences:
  they raise output and the domestic interest rate, and may trigger retaliatory
  tariffs, trade wars, or exchange rate adjustments that partly or fully undo
  the initial stimulus.
\end{solution}
\end{question}

\begin{question}[type=open]{Why Should Central Banks Be Independent of Government,
  and What Tensions Does This Create?}
The script notes that ``most central banks today act fully independent of
the government.'' Explain the economic rationale for central bank independence
and the potential conflicts this creates with elected governments.

\answerlines[7]

\begin{solution}
  \textbf{Rationale for independence:} Governments face electoral incentives
  to stimulate the economy before elections, even if doing so causes
  inflation. If the central bank were under government control, monetary
  policy would be systematically too loose, leading to a chronic inflation
  bias (Kydland-Prescott, time-inconsistency problem). An independent central
  bank can credibly commit to price stability because it does not face
  re-election pressure.

  \textbf{The IS-LM tension:} Fiscal policy (IS curve) is set by the
  government; monetary policy (LM curve) is set by the independent central
  bank. If the government runs a large fiscal stimulus (rightward IS shift),
  the central bank may counter with tighter money (leftward LM shift) to
  contain inflation. The result is higher interest rates but little change in
  output --- the two policies partially cancel out. This is the
  ``obvious contradiction'' the script mentions between fiscal and monetary
  policy.

  \textbf{Crisis tensions:} During recessions, governments want the central
  bank to finance deficits (``monetary financing''). An independent central
  bank can refuse. The Eurosystem's prohibition on direct government financing
  (Art.\ 123 TFEU) was heavily contested during the Eurozone debt crisis and
  again during COVID-19.
\end{solution}
\end{question}

\end{document}